\appendix
\onecolumngrid

\section{Notation}
\label{app:notation}

Table~\ref{tab:symbols} defines every symbol of this paper. Symbols marked ``taken'' are those of the
rest of PROJECT spacetimes, used here with exactly the meaning that they have there (their meaning is
copied, with a note on their use here where needed); all others are new. Every new symbol of the gauge
sector and of the superfluid carries the suffix U2, so that no letter of this paper can be read as a
symbol of the same letter in another part of the project; the Dirac matrices and their relatives keep
the tag 1671 of the $7+1$ spacetime, and the symbols of the electroweak theory, quoted as the template,
carry the suffix SM. The machine checks R3-S1 to R3-S6 (Appendix~\ref{app:checks}) verify that no
symbol is registered twice, that the taken symbols carry the meaning of the project's master
registries, that no new symbol coincides with a symbol of those registries, and that no letter with
another meaning in the project is used in a formula of this paper without its suffix. Units: the
cutoff $\Lambda_{U2}=1$ unless a unit is shown; $\hbar=1$ and the speed of light is 1.

{\small
\begin{longtable}{p{0.24\textwidth}p{0.54\textwidth}>{\raggedright\arraybackslash\footnotesize}p{0.16\textwidth}}
\caption{The symbols of this paper, by group, with the names under which the programs and the data
files of this paper write them.}\label{tab:symbols}\\
\toprule
symbol & meaning & names in data files\\
\midrule
\endfirsthead
\toprule
symbol & meaning (continued) & names in data files\\
\midrule
\endhead
\bottomrule
\endfoot
{{GENRAW:symbols_r3_texrows}}
\end{longtable}
}

\section{The machine checks}
\label{app:checks}

The program \texttt{dark\_u2} (Appendix~\ref{app:data}) verifies {{GEN:nchecks_rust}} statements on
which the results rest, the script \texttt{check\_symbols.py} the six statements on the symbols, and
\texttt{check\_self\_contained.py} that no document of this revision sends the reader to another
document for a definition. Each check prints one line: its label, the statement, and, after two
dashes, the numbers that decide it. Table~\ref{tab:checks} reproduces these lines as the programs
print them, in plain text: \texttt{gamma\^{}A\_1671} is $\gamma^A_{1671}$, \texttt{eta1671} is
$\eta_{1671}$, \texttt{omega} is $\omega_{1671}$, \texttt{C\^{}+-} is $C^\pm_{1671}$, \texttt{beta},
\texttt{alpha\^{}j} are $\beta_{1671}$, $\alpha^j_{1671}$, \texttt{Psi} is $\Psi_{U2}$, \texttt{tau\^{}c}
is one of the Pauli matrices $\tau^1_{U2}$, $\tau^2_{U2}$, $\tau^3_{U2}$, \texttt{Ups} is
$\Upsilon_{U2}$, \texttt{theta}, \texttt{theta0} are $\boldsymbol\Theta_{U2}$, $\Theta^0_{U2}$,
\texttt{calW}, \texttt{calF} are $\mathcal W_{U2A}$, $\mathcal F_{U2AC}$, \texttt{g}, \texttt{g'},
\texttt{y} are $g_{U2}$, $g'_{U2}$, $y_{U2}$, \texttt{Delta} is $\Delta_{U2}$, \texttt{mu} is
$\mu_{U2}$, \texttt{M} and \texttt{m} are $M_{U2}$ and $m_{U2}$, \texttt{eps} is $\varepsilon_{U2}$,
\texttt{E\^{}+-} is $E^\pm_{U2}$, \texttt{xi}, \texttt{u}, \texttt{v} are $\xi^\pm_{U2}$, $u^\pm_{U2}$,
$v^\pm_{U2}$, \texttt{h(q)} is $\mathsf h_{U2}(\vec q)$, \texttt{w} in R3-C3 is an eigenvector
$\mathsf e^\pm_{U2}$ and elsewhere the frequency $w_{U2}$, \texttt{q\_F} is $q^F_{U2}$, \texttt{Q} is
$\vec Q_{U2}$, \texttt{G\^{}P}, \texttt{G\^{}S}, \texttt{G\^{}P\_crit} are $G^P_{U2}$, $G^S_{U2}$,
$G^{\rm crit}_{U2}$, \texttt{lambda\^{}P}, \texttt{lambda\^{}S} are $\lambda^P_{U2}$, $\lambda^S_{U2}$,
\texttt{a7}, \texttt{b8} are $a^{(7)}_{U2}$, $a^{(8)}_{U2}$, \texttt{F(s)} is $L_{U2}(s_{U2})$,
\texttt{J\_n} is $\mathcal J^{(\iota)}_{U2}$ with $\iota$ the number written in place of n,
\texttt{Omega} is $\Omega_{U2}$, \texttt{n} is $n_{U2}$, \texttt{f\_s}, \texttt{f\_t}, \texttt{f}
are $f^s_{U2}$, $f^t_{U2}$, $f_{U2}$, \texttt{m\_sigma} is $m_{\sigma U2}$, \texttt{U''} is the diagonal
matrix of Eq.~\eqref{eq:R}, \texttt{Pi} is $\Pi_{U2}$, \texttt{T\_sigma} is $\mathcal V^\sigma_{U2}$,
\texttt{m\_B} is $m_{BU2}$, \texttt{G7(r)} is $V^{\rm C}_{U2}(\varrho_{U2})$ and \texttt{phi} is
$\phi^{\rm Y}_{U2}$, \texttt{T\_1.T\_2} is $\tfrac14\boldsymbol\tau^{(1)}_{U2}\cdot\boldsymbol\tau^{(2)}_{U2}$,
\texttt{Y} in R3-B2 is twice the charge, and \texttt{x} in R3-G7 is $m_{BU2}\varrho_{U2}$. In
R3-A6, R3-A12 and R3-B2 the letters \texttt{S} and \texttt{K} are a generator and a bilinear form,
$\Sigma^{AC}_{1671}$ (up to the factor $\tfrac14$) and $\mathsf K$.

{\footnotesize
\begin{longtable}{p{0.08\textwidth}p{0.80\textwidth}p{0.06\textwidth}}
\caption{All machine checks of this paper, as the programs print them.}\label{tab:checks}\\
\toprule
label & statement and decisive numbers & result\\
\midrule
\endfirsthead
\toprule
label & statement and decisive numbers (continued) & result\\
\midrule
\endhead
\bottomrule
\endfoot
{{GENRAW:checks_rows_tex}}
\end{longtable}
}

\section{Programs, data, and how to reproduce every number}
\label{app:data}

\subsection{Programs}

All computations are done by the Rust program \texttt{dark\_u2} in the folder
\texttt{code/dark-u2-rs} of \texttt{Revision\_3}; the figures are drawn by a Python script, and the
documents are assembled by Python scripts that insert every number from the data files, so that no
computed number is typed by hand. The program has nine modules: \texttt{algebra} (the Dirac matrices,
the pairs, the doublet, the anomaly coefficients; exact integer and Gaussian-integer arithmetic, ranks
modulo the prime $2^{61}-1$), \texttt{gauge} (covariance of $D_{U2A}$ with random fields at a point,
derivatives by fourth-order central differences; the static potentials), \texttt{bdg} (the
Bogoliubov--de Gennes matrix, its spectrum and spinors), \texttt{potential} (the effective potential,
the gap and mass equations, the minima), \texttt{modes} (the random-phase approximation, the
stiffness), \texttt{integrate} (adaptive Gauss--Legendre quadrature with 20 and 40 nodes, relative
tolerance $10^{-15}$), \texttt{rsym} (real symmetric matrices and their cyclic Jacobi
diagonalization), \texttt{checks} (the reports) and \texttt{main}. It uses the library crate
\texttt{code/spacetimes-rs} of the repository for complex numbers, integer matrices and their rank,
Gauss--Legendre nodes, Brent's root finder, a parallel map and the writing of data files.

Numerical methods. Near a Fermi surface the integrals over $q$ use the variable
$\mathrm{arsinh}(\xi^+_{U2}/|\Delta_{U2}|)$, in which the integrands are smooth for every gap however
small; $E^-_{U2}$ is then computed as $|\Delta_{U2}|\cosh$ of that variable, without rounding error. The
gap equation is solved by Brent's method in $\ln|\Delta_{U2}|$ (gaps down to $10^{-300}$); the global
minimum over $M_{U2}$ and $\Delta_{U2}$ is found on the profile $\min_{M_{U2}}\Omega_{U2}$ with the gap
solved at each $M_{U2}$, by golden-section search and a final solution of the mass equation. The
polarization uses Gauss--Legendre panels refined around the Fermi momentum, with $\vec q$ along the
coordinate $x1$ (the integrands depend on the length of $\vec q$ only), and the eigenvectors of the
$32\times32$ matrix at each node. The mass of $\sigma_{U2}$ is the zero of $\det\mathsf R_{U2}(w_{U2})$
below the threshold, or the threshold itself.

\subsection{Commands}

From the folder \texttt{Revision\_3} of the repository (a Rust toolchain for the edition 2021,
Python 3 with NumPy and Matplotlib, and for the documents a TeX distribution with REVTeX 4.2, pandoc
and XeLaTeX):
\begin{verbatim}
cargo run --release --manifest-path code/dark-u2-rs/Cargo.toml -- all
python3 code/scripts/make_figures.py
bash code/scripts/check_all.sh
\end{verbatim}
The first command writes every data file of the folder \texttt{data} (in a few seconds); the commands \texttt{algebra}, \texttt{gauge}, \texttt{bdg}, \texttt{potential} and
\texttt{modes} in place of \texttt{all} run one part. The second draws the nine figures. The third
runs both, then the checks of the symbols and of self-containedness, writes the fragments of the
documents, and builds this paper and the provenance documents.

\subsection{Data files}

\texttt{r3\_algebra\_checks.txt}, \texttt{r3\_gauge\_checks.txt}, \texttt{r3\_bdg\_checks.txt},
\texttt{r3\_potential\_checks.txt}, \texttt{r3\_modes\_checks.txt}, \texttt{r3\_symbol\_checks.txt} and
\texttt{r3\_self\_contained\_checks.txt}: the reports of Table~\ref{tab:checks}.
\texttt{r3\_gamma\_kron.csv}: the Kronecker factors of Eq.~\eqref{eq:gammas}.
\texttt{r3\_pairs\_by\_rank.csv}: Table~\ref{tab:pairs}.
\texttt{r3\_static\_potential.csv}: Fig.~\ref{fig:static}.
\texttt{r3\_fermi\_surface\_gap.csv} and \texttt{r3\_spectrum.csv}: Fig.~\ref{fig:spectrum}.
\texttt{r3\_hat\_radial.csv} and \texttt{r3\_hat\_2d.csv}: Fig.~\ref{fig:hat}.
\texttt{r3\_vacuum\_gap.csv}: Fig.~\ref{fig:gap} and Table~\ref{tab:vacuum}.
\texttt{r3\_competition.csv}: Fig.~\ref{fig:competition}.
\texttt{r3\_density\_scan.csv} and \texttt{r3\_bcs\_slope.csv}: Fig.~\ref{fig:density}.
\texttt{r3\_phase\_map.csv}: Fig.~\ref{fig:phase}.
\texttt{r3\_transition.csv}: Fig.~\ref{fig:transition}.
\texttt{r3\_modes\_vacuum.csv}, \texttt{r3\_modes\_mass.csv} and \texttt{r3\_modes\_density.csv}:
Fig.~\ref{fig:modes} and Tables~\ref{tab:sigma} to~\ref{tab:density}.
\texttt{r3\_summary\_potential.csv} and \texttt{r3\_summary\_modes.csv}: the key numbers quoted in the
text. Every data file states its meaning and units in its header lines, and every column name is
either a descriptive word or one of the names listed in the last column of Table~\ref{tab:symbols}.

\subsection{Checksums}

SHA-256 checksums of the data files, the figures and the programs from which this paper was built:
{\footnotesize
\begin{verbatim}
{{GENRAW:data_sha256}}
{{GENRAW:figure_sha256}}
{{GENRAW:code_sha256}}
\end{verbatim}
}

\twocolumngrid
